\makeatletter \def\input@path{{../../}} \makeatother
\documentclass{article}
\usepackage[utf8]{inputenc}
\usepackage{amsmath, amsthm, amssymb, mathpazo, isomath, mathtools}
\usepackage{subcaption,graphicx,pgfplots}
\usepackage{fullpage}
\usepackage{booktabs}
\usepackage{hyperref}
\usepackage{algorithm, algorithmic}
\usepackage{mathtools}
\usepackage{todonotes}
\usepackage{tcolorbox}
\usepackage{totcount,xfp}
\newtotcounter{totalpts}
%\setcounter{totalpts}{0} % Contador de puntos en homework

\newcommand{\getTotal}{\number\totvalue{totalpts} }
\newcommand{\instruccion}{\begin{tcolorbox}[colback=blue!5!white,colframe=blue!75!black!70!white,title=Instructions]
The homework must be developed by one or two people in a .pdf report submitted through the designated drop box on the Canvas platform, where you must also show the developed code if applicable. For your convenience, you may submit the homeworks in a Jupyter Notebook, so it is more comfortable to show the code. In any case, a single file must be submitted as the answer to the homework. The late policy will be: a linear factor will be calculated that is 1 at the time of delivery and 0 48 hours later. This will multiply your obtained score.


The homework has \getTotal points. In each question, the given score will be: the total if it is correct, half if it has a conclusive development but contains errors, and 0.0 points if it has little to no development.

The use of language models (like ChatGPT) is allowed to support the development of the homework, although we suggest trying each question a bit first to generate doubts. While we cannot control if you ask the model directly for the answer, we believe it is better for your learning process to ask for suggestions rather than solutions. For this reason, if you decide to use it, you must attach to your homework another pdf showing the prompts used. We suggest using prompts of this style: "I am a student of [Major] and I am developing the following exercise for a homework in the course [Course]: [Copy exercise]
  And we have seen so far the contents in the attached slides [attach pdfs of the lectures]. Given this content, I would like to ask you for a suggestion to solve [explain difficulty], since I am confused about [write doubt]. As an answer, I do not want you to give me the solution, but a guide to help me start developing my answer."

  If you have doubts and/or comments, please contact us by email or Canvas. It is possible that the formulation has errors.
\end{tcolorbox}}


\newcommand{\code}[1]{\texttt{#1}}
% Logic: To create a counter with the points in each question and accumulate that in a total variable. Now I simply set the full threshold at the beginning to 75% and it gives the computed value automatically.
\newcommand{\pts}[1]{\if#11\textbf{[#1 point]}\else\textbf{[#1 points]}\fi\addtocounter{totalpts}{#1}}


\title{Homework 1}
%\author{Nicol\'as A Barnafi\thanks{Instituto de Ingeniería Biológica y Médica, Pontificia Universidad Católica de Chile, Chile}, Axel Osses\thanks{Departamento de Ingeniería Matemática, Universidad de Chile, Chile}}
%\author{Nicol\'as A Barnafi}
\date{}

\renewcommand{\vec}{\vectorsym}
\newcommand{\mat}{\matrixsym}
\newcommand{\ten}{\tensorsym}
\newcommand{\der}[2]{\frac{d#1}{d#2}}
\DeclareMathOperator{\grad}{\nabla}
\DeclareMathOperator{\dive}{\text{div}}
\DeclareMathOperator{\curl}{\text{curl}}
\DeclareMathOperator{\tr}{\text{tr}}
\DeclareMathOperator{\sym}{\text{sym}}
\newtheorem{remark}{Remark}
\newtheorem{definition}{Definition}
\newcommand{\R}{\mathbb{R}}
\newcommand{\D}{\mathcal{D}}
\newcommand{\parder}[2]{\frac{\partial\,#1}{\partial\,#2}}

\newcommand{\tin}{\text{in}}
\newcommand{\ton}{\text{on}}

\newtheorem{theorem}{Theorem}
\newtheorem{lemma}{Lemma}

\begin{document}

\maketitle

\instruccion

\begin{enumerate}
    \item\pts{2} Define the higher-order Fréchet derivatives. Then, consider the operator $T: L^p(\Omega) \to \R$ given by $T(f) = \int_\Omega f^p\,dx$ and compute its third-order derivative.
    \item\pts{3} State and prove the chain rule for Fréchet derivatives, given by $d(G\circ F)(u)[h] = dG(F(u))[dF(u)[h]]$. Then, use it to find the derivative of the composition $S\circ T$ of the operators $S:L^\infty(\Omega)\to L^\infty(\Omega)$ and $T:L^\infty(\Omega)\to L^\infty(\Omega)$ given by $S(u) = \sin u$ and $T(u) = \cos u$. Precisely define all the spaces used to apply the chain rule. 
    \item\pts{2} The Cayley-Hamilton Theorem states that every matrix satisfies its own characteristic equation using the invariants: 
                $$ \mat A^3 - I_1(\mat A) \mat A^2 + I_2(\mat A) \mat A - I_3(\mat A) \mat I = 0. $$
                Differentiate this equation with respect to $\mat A$ and expand the equation to prove that $\parder{\det \mat A}{\mat A} = \det (\mat A) \mat A^{-T}$.
    \item\pts{1}  Use the previous question to prove that if $A:\R \to \R^{n\times n}$ is a matrix function, then
            $$ \der{\det A}{t}(t) = \det A(t) A^{-T}(t) : \der A{t}(t). $$
            \item If $\phi, \vec u, \ten A$ are a scalar, vector, and tensor field respectively in $\R^3$, prove using index notation that if $\vec v$ is another vector field:
                    \begin{itemize}
                        \item\pts{1} $$\dive(\phi\vec  u) = \phi \dive \vec u + \vec u \cdot \grad \phi $$
                        \item\pts{1} $$\dive(\phi\ten  A) = \phi \dive \ten A + \ten A \grad \phi $$
                        \item\pts{1} $$\dive(\vec u \times \vec v) = \vec v\cdot \curl \vec u - \vec u \cdot \curl v $$
                        \item\pts{1} $$ \grad(\phi \vec u) = \vec u \otimes \grad \phi + \phi \grad \vec u$$
                    \end{itemize}
    \item Consider the domain $\Omega=(a,b)$ with $a<b$.
            \begin{itemize}
                \item\pts{1} Consider the function $f:\Omega \to \R$ given by $f(x) = \sin (2\pi x)$. For it, consider a numerical differentiation scheme using the finite difference method to approximate its derivative $f'(x)$ in $\Omega$. In particular, compute the theoretical convergence rate for the forward, backward, and centered difference schemes in the $\|\cdot \|_{\infty,h}$ norm. For this, consider $a=-1$, $b=1$, and a distance between discrete points given by $h  \in \{2^{-n}\}_{n=4}^{20}$, compute the empirical convergence rates and interpret the obtained results.
                \item\pts{1} Repeat the previous exercise for the function $f(x) = 2$ and interpret the obtained results.
                \item\pts{1} For this problem, when computing the numerical derivatives at the boundaries, you will have issues when using first-order finite differences combined with centered differences, which will deteriorate convergence. While this is not a problem, you can choose to use a one-sided finite difference formula to obtain quadratic convergence at the endpoints. By computing the Taylor series associated with $f(x+h)$ and $f(x+2h)$, prove that the corresponding approximation is
            $$ f'(x) = \frac{-3f(x) + 4 f(x+h) - f(x+2h)}{2h} + O(h^2). $$
            Use this formula to show that the quadratic convergence of centered differences is recovered.
            \end{itemize}
        \item\pts{3} Implement the first-order discrete derivatives $\ten D^+, \ten D^-, \ten D$ and construct an example to validate the convergence rates seen in class. Analytically compute the convergence rate of centered differences for the second-order derivatives $\ten D^-\ten D^+$, and program an example to also validate its convergence rate. Finally, analyze the case of the composition of centered differences $\ten D\ten D$.
        \item\pts{2} The notes detail how centered differences can be written as the composition of the forward and backward difference operators according to $\ten D^+ \ten D^- = \ten D^-\ten D^+$. Prove theoretically that instead $\ten D^+\ten D^+$ (or $\ten D^-\ten D^-$) is not a good approximation of the second-order derivative.
    \item Consider the domain $\Omega=(0,1)$. In it, discretize the equation 
            $$ \mathcal L u = 1 $$ 
            with the differential operator $\mathcal L$ given by
            $$ \mathcal L u = - u'' + bu' + cu, $$
            with boundary conditions given by $u(0) = 1$, $u(1) = 2$. This problem is known as the Advection-Diffusion-Reaction (ADR) equation. We will compute the sensitivity of the problem to its parameters in terms of the conditioning of the system, which is intrinsically related to the numerical difficulty of the problem. We will call the matrix induced by the discretization $\mat L_h$. 
            \begin{itemize}
                \item\pts{1} Show the corresponding equation for the grid node $i$ for (i) forward differences for $u'$ and (ii) centered differences for $u'$. 
                \item\pts{2} In each of the two discretizations from the previous point, show that the obtained discretization is consistent and compute the order of consistency.
                \item\pts{1} Plot in a log-log graph the conditioning of the matrix $\mat L_h$ with respect to the following values: $h \in \{2^{-n}\}_{n=2}^8$, $b \in \{0.01,0.1,1,10,100\}$, and $c \in \{0.01,0.1,1,10,100\}$. Do this only for the scheme where $u'$ is approximated with centered differences.
                \item\pts{2} Prove theoretically using the Lax Equivalence Theorem that the discrete solution obtained with the matrix $\mat L_h$ converges to the continuous solution if $b=0$. To do this, extend the proof from the notes for the Laplacian.
                \item\pts{2} Using the method of manufactured solutions, show that the numerical convergence rate to an analytical solution is the one expected by the theory. 
            \end{itemize}

\end{enumerate}

\todo[inline,color=white!90!black]{\textbf{Note: } We will open a forum on Canvas to review any typos and/or errors in the formulation.}
\end{document}
